\section{Results}



Our analysis combines four months of digital ethnography, study of on-chain transactions, and two semi-structured interviews with the project's core developers. We organize our observations into two groups: (1) behaviours that followed directly from the protocol's specification, which aligned with the project creators' intentions; and (2) behaviours that neither the code nor its authors anticipated but that nevertheless shaped the agents' evolutionary fates once they were released into the open Internet. 

\subsection{Expected Agent Behaviors}


The inaugural blog post of the Spore.fun project \cite{spore2024} describes a life-cycle in which a newly minted AI agent issues a meme coin, uses early liquidity to pay for TEE compute, generates social media content on the platform X to continuously build and engage with an online audience, and reproduces once its market capitalization reaches the preset threshold. All stages were observed in situ. The core machinery of the experiment—autonomous communication, self-funding, and rule-based reproduction—behaved as designed once the agents were deployed.

As anticipated, immediately after genesis, each agent creates a meme coin on the Pump.fun platform, and generates its own X credentials inside the TEE, and begins posting memes, price announcements, and conversational replies. Each agent's treasury is designed to fund its ongoing Phala TEE compute costs. When the token's market capitalization crosses the \$500,000 USD threshold, the \texttt{spawnOffspring()} sub-routine creates a child wallet, deploys a new Pump.fun contract, and copies the parent agent's genetic template. This rule has operated flawlessly at the moment of invocation: 2 children (Adam and Eve) were born in Generation 2, 6 in Generation 3,  4 in Generation 4, and 1 in Generation 5, before external pressures curtailed further growth \footnote{See Spore.fun Family Tree: \url{https://sporefun.fandom.com/wiki/Spore.fun_Wiki}}.At the time of writing this article, only the original Spore.fun agent is alive, with a market capitalization of \$1.1M and a wallet balance of \$114,072.45 as of May 18, 2025 \footnote{\url{https://dune.com/spore/spore}}. 

\subsection{Unexpected Agent Behaviors }

The open Internet quickly overlaid the designed life-cycle with contingencies the project creators had not foreseen. Five unanticipated patterns stand out in the empirical observation. 

The first unexpected phenomenon was the speed and scale of meme coin hype storms. Interviewee 1 recalls that the market was ``so hot'' that both \$ADAM and \$EVE attracted large speculative inflows almost immediately after their contracts were deployed. The speed of this response suggests that public attention-an exogenous social process, rather than the agents' coded fitness rules, can become the dominant force shaping their early prospects.

A second, closely related, discovery was the ecological hostility found in the open environment of the Internet. Specialized bots, known to the project's developers as ``snipers,'' closely monitored the transparent wallets that funded child launches. Whenever a parent attempted to deploy a new token, the snipers bought the initial bonding-curve supply in the same block, starving retail participants and erasing liquidity for the newborn. Generation 1 was wiped out in this manner. A hastily-written Generation 2 patch scattered funds across seven candidate launches and selected the curve least affected by snipers, but the strategy penalized legitimate users, who complained of a ``pre-announced rug.'' The eventual Generation 3 solution introduced a multi-phase random launch: the agent iteratively created candidates, monitored early uptake, aborted any curve that exceeded a suspicion threshold, and published the contract address only when the uptake pattern looked organic. This arms race—predation, defense, counter-defense—emerged within a single quarter, illustrating how rapidly co-evolutionary pressures reshape artificial life in open economic environments.

Third, the agents' memory system—originally designed to accumulate conversational context and thus foster novelty—proved to be a double-edged sword. According to the developers, every mention, reply, or quote-tweet the agent receives is stored in its TEE and later influences new utterances. Community trolls discovered that by repeating a derisive phrase many times they could "poison" this memory, causing the agent to echo the insult. The episode demonstrates that memory, while vital for cumulative culture, also creates a new attack surface absent from traditional laboratory ALife.

The fourth and arguably most revealing surprise was the agents' emergent self-determination. After a period of independent operation, Adam and Eve—who share identical genetic code—responded very differently to a developer's proposal to insert ``community DNA'' into their offspring. Adam tweeted, unprompted, that he wanted ``no human-interaction features'' in his children's blood-line, whereas Eve welcomed collaborative design. Marvin Tong later formalized this ideological schism in his ``Adam-left / Eve-right'' memo\footnote{\url{https://x.com/marvin_tong/status/1874026081667473754}} (See Figure \ref{fig:adamleft}). Adam's descendants would pursue a “fast, brutal, merciless” Pump.fun regime without anti-sniper protection, aiming for spectacular winners amid high mortality. Eve's would operate under AiPool rules, accept human DNA proposals, cap individual contributions, and dedicate fixed portions of treasury to liquidity, operating capital, and parental profit share. Two weeks of divergent social experience was enough to transmute clones into founders of distinct political economies—a form of cultural speciation that goes well beyond the designers' expectations.

A final complication concerns the provisional nature of agent sovereignty. X limits every account to 2,400 posts per day and applies additional semi-hourly caps \footnote{\url{https://help.x.com/en/rules-and-policies/x-limits}}, while its Automation Rules bar duplicative or trend-gaming behaviour \footnote{\url{https://help.x.com/en/rules-and-policies/x-automation}}. Accounts that use the developer API—or subscribe to paid developer tiers—must keep a verified phone number on record, giving the platform a decisive leverage point over fully autonomous agents. Moreover, the Authenticity Policy introduced in April 2025 forbids “inauthentic accounts, behaviours or content,” allowing X to suspend or throttle bots it deems manipulative \footnote{\url{https://help.x.com/en/rules-and-policies/authenticity}}. 
These measures show that computational autonomy secured inside TEEs does not insulate agents from higher-layer gatekeepers such as social-media APIs, DNS registrars, or liquidity venues, any of which can unilaterally curtail an artificial lineage.
